\documentclass[11pt]{article}

% bioRxiv-compatible packages
\usepackage[utf8]{inputenc}
\usepackage[T1]{fontenc}
\usepackage{amsmath,amssymb}
\usepackage{graphicx}
\usepackage[margin=1in]{geometry}
\usepackage{natbib}
\usepackage{booktabs}
\usepackage{hyperref}
\usepackage{xcolor}
\usepackage{setspace}
\usepackage{textcomp}

\onehalfspacing

% Title
\title{Extensible Longevity Model (ELM): A Calibrated Interventions Framework for Aging Research}
\author{Silmon James Biggs}
\date{}

\begin{document}

\maketitle

\begin{abstract}
We present a mechanistic ODE model of mammalian aging with 6 calibrated parameters, fitted to male ITP lifespan data by root-finding. The model predicts female ITP outcomes for all six compounds with no sex-specific fitting, using four sex mechanisms drawn from prior literature, achieving 1.3 percentage point mean error. To our knowledge, this is the first quantitative model that predicts the direction and magnitude of ITP sex differences across multiple compounds from male data alone. The model also predicts the rapamycin + acarbose combination outcome (+33.5\% predicted vs.\ +34\% observed) without fitting to that result. Four coupled aging subsystems---mitochondrial heteroplasmy, DNA damage, epigenetic drift, and cellular senescence---are linked through NAD$^+$ metabolism and SASP signaling. The dominant downstream prediction is a heteroplasmy bottleneck: after the best ITP compound stack, mitochondrial heteroplasmy is the largest unaddressed contributor to residual biological age, identifying mitophagy-enhancing compounds as the highest-value novel targets.
\end{abstract}

% ============================================================
\section{Introduction}
% ============================================================

The biology of aging has attracted sustained theoretical attention for more than a century, generating a literature remarkable more for its breadth than for its predictive power. More than 300 distinct theories of aging have been catalogued \citep{Lemoine2021}, spanning evolutionary arguments grounded in life-history trade-offs \citep{Medawar1952,Williams1957,Kirkwood1977}, thermodynamic accounts invoking entropy \citep{Hayflick2007}, programmatic hypotheses, and network theories of damage accumulation \citep{Kowald1997}. This body of work addresses a genuine and important question: why do organisms age? It has been valuable in organizing biological intuition and motivating experimental inquiry.

The present work addresses a different question: given that aging happens, what do specific interventions do to its rate, and which combinations are worth testing next? The distinction---between a theory of aging causation and a quantitative model of aging interventions---is the difference between explaining why an organism dies and predicting when. A causation theory is evaluated on whether it makes sense; an interventions model is evaluated on out-of-sample prediction. We describe the Extensible Longevity Model (ELM), a mechanistic ordinary differential equation model of the aging process, calibrated to the best available mammalian lifespan data and designed to predict the outcomes of multi-compound intervention strategies.

\subsection{The Data Foundation: NIA Interventions Testing Program}

Any quantitative interventions model requires a quantitative target. The NIA Interventions Testing Program (ITP) provides the most rigorous mammalian lifespan dataset currently available \citep{Miller2007,Strong2008,Harrison2009}. The ITP tests candidate compounds in genetically heterogeneous UM-HET3 mice simultaneously across three independent sites, controlling for laboratory-specific artifacts that have historically limited the reproducibility of single-site longevity studies \citep{Nadon2017}. The program has identified six compounds with reproducible, statistically robust lifespan extension: rapamycin \citep{Miller2011rapa}, acarbose \citep{Harrison2014}, 17$\alpha$-estradiol \citep{Stout2017}, canagliflozin \citep{Miller2020}, aspirin \citep{Strong2016}, and glycine \citep{Miller2019glycine}. These compounds span diverse mechanisms---mTORC1 inhibition, glycemic modulation, antioxidant activity, anti-inflammatory signaling---and show striking sex differences, with four of six producing male-biased or male-only effects.

Each compound has been tested individually. The question of which combinations are worth pursuing---and why---remains unanswered, because the combinatorial space is too large to screen exhaustively and no quantitative model for predicting combination outcomes from single-compound data has been available. The ELM is built to fill that gap, using the ITP's single-compound results as its calibration targets and the biological mechanisms underlying those results as its structural skeleton.

\subsection{Why Single Compounds Hit a Ceiling: Reliability Theory}

Before describing the model, it is worth establishing why the ITP's best single-compound results---lifespan extensions of +6\% to +26\%---should be understood not as a biological ceiling but as a mathematical consequence of experimental design.

Gavrilov and Gavrilova's reliability theory of aging \citep{Gavrilov2001,Gavrilov2006} treats organisms as series systems: structures that fail when any one critical subsystem reaches its failure threshold. The organism dies from whichever mode strikes first. The consequence is that fixing one failure mode while leaving others intact produces gains proportional to that mode's contribution and no more. The remaining modes do not vanish; they wait. Under equal-weight exponential failure assumptions, $n$ independent failure modes yield a maximum single-mode extension of $\sim 1/(n-1)$; the observed +10--26\% range from the best single ITP compounds is consistent with $n \approx 4$--10.

The most informative experiments are those that address multiple failure modes at once, because single-compound testing---however rigorously conducted---cannot see what is behind a wall it has not yet removed. The ELM is designed to map those modes, their interactions, and the combinations most likely to address more than one simultaneously.

\subsection{Biological Architecture: The Hallmarks as Quantitative Targets}

The choice of which aging processes to model is constrained by two requirements: the processes must respond to at least one ITP compound (providing calibration leverage), and their dynamics must be expressible in mechanistically grounded differential equations with literature-constrained parameters. Four processes satisfy both criteria.

These correspond to a subset of the Hallmarks of Aging framework introduced by L\'opez-Ot\'in and colleagues \citep{LopezOtin2013,LopezOtin2023}: mitochondrial dysfunction (operationalized as mitochondrial heteroplasmy), genomic instability (DNA damage accumulation), epigenetic alterations (methylation drift), and cellular senescence. The hallmarks framework provides the biological vocabulary; the ELM attempts to put calibrated numbers on the subset of it that is experimentally tractable with current ITP data.

The four subsystems are not independent. They are coupled through shared molecular intermediates, most centrally NAD$^+$ \citep{Yoshino2011,CamachoPereira2016}, whose availability governs DNA repair (PARP1 \citep{Bai2011}), mitophagy (SIRT1/FOXO3 \citep{Imai2014,Warr2013}), and epigenetic stability (EZH2 \citep{Vire2006}). A drug that raises NAD$^+$ therefore has the potential to partially address multiple hallmarks simultaneously; conversely, two drugs that both raise NAD$^+$ may show diminishing returns when combined. Senescent cells add a fifth player: they secrete SASP \citep{Coppe2008}, which upregulates the NAD$^+$-consuming ectoenzyme CD38 \citep{Tarrago2018} and accelerates oxidative damage, creating a feedback loop between cellular senescence and NAD$^+$ decline. The ODE system (Section~2) captures these couplings explicitly.

\subsection{Sex Differences: Mechanisms from Prior Literature}

The ITP sex differences are among the most striking and underexplained features of the dataset. Jiang et al.\ \citep{Jiang2025} catalogued sex-specific outcomes across 54 compounds and 164 ITP trials over 20 years, documenting the pattern extensively. Garratt \citep{Garratt2020} reviewed qualitative hypotheses about sex-specific pathways, including growth hormone/IGF-1 and mTOR signaling. But no prior work has provided a quantitative model that predicts the direction and magnitude of ITP sex differences across multiple compounds from male calibration data alone. The individual sex mechanisms we employ are each drawn from prior experimental literature; their unification into a single predictive model, tested against ITP data it was not fitted to, constitutes the primary contribution of this paper.

The four mechanisms are: (i)~AMPK saturation---females have higher baseline AMPK via LKB1 upregulation \citep{Boudaba2018,McInnes2011}, attenuating response to AMPK activators; (ii)~CYP3A pharmacokinetics---females achieve higher rapamycin blood levels \citep{Lamming2014}; (iii)~testosterone dependency---17$\alpha$-estradiol requires 5$\alpha$-reductase substrate \citep{Garratt2018}; (iv)~COX-2/prostacyclin interference---aspirin suppresses estrogen-mediated cardioprotection \citep{Egan2004,Grosser2006}. Each is drawn from prior experimental literature, not fitted to ITP sex data. Whether these assignments correctly predict the observed sex ratios is evaluated in the Results.

\subsection{Scope and Limitations}

The ELM is a model of metabolic aging in tissues that depend heavily on oxidative phosphorylation. It does not model Wnt/$\beta$-catenin or Notch signaling, telomere dynamics, thymic involution, immune senescence, AGE crosslink accumulation, or organ-specific fibrosis. These processes may become rate-limiting at substantially extended lifespans and would require separate modeling frameworks or extension of the ELM.

Three recently described rejuvenation modalities---partial cellular reprogramming via OSK factors \citep{Ocampo2016}, EZH2 non-catalytic rejuvenation \citep{Liao2024}, and telomere transfer via extracellular vesicles from CD4$^+$ T cells \citep{Lanna2025}---lie outside the model's current scope. The last of these achieved +65\% median lifespan extension in aged C57BL/6 mice in a preprint from Lanna, Valvo, and Dustin \citep{Lanna2025}, suggesting that the biological ceiling for intervention is substantially higher than single-compound ITP results imply, and that multiple mechanistically independent axes of aging remain untapped.

The model's parameters are of three types: those calibrated to ITP lifespan data, those drawn from published experimental literature, and those estimated from biological coupling structure. The model makes no claim to completeness---it is, as its name indicates, designed to be extended as new data become available. The ELM is a calibrated prior for which intervention targets matter most. Its predictions are falsifiable, its assumptions are explicit, and its purpose is to inform decisions about which combinations may warrant testing.

% ============================================================
\section{Methods}
% ============================================================

\subsection{Model Structure}

The ELM is a system of coupled ordinary differential equations describing the dynamics of four aging subsystems over the lifespan of a mouse. A scalar BioAge composite integrates the four subsystem states into a single mortality-predictive variable, which is mapped to survival probability through a Gompertz hazard function. Compound effects enter as perturbations to specific rate constants, with entry points and magnitudes derived from published mechanistic and pharmacological literature. Code is available at \url{https://github.com/silmonbiggs/extensible-longevity-model}.

\subsubsection{State variables}

Four primary state variables track the aging subsystems. Mitochondrial heteroplasmy ($H$) is the fractional abundance of mutant mtDNA; initial value $H(0) = 0.011$, bounded $[0, 1]$. NAD$^+$ level ($N$) is normalized intracellular concentration relative to young-adult baseline; $N(0) = 1.0$. DNA damage ($D$) is normalized cumulative double-strand break equivalents; $D(0) = 1.0$. Senescent cell burden ($S$) is normalized p16$^{\text{Ink4a}}$-positive cell abundance; $S(0) = 1.0$. A fifth variable methylation drift ($M$), initialized at $M(0) = 0$, tracks deviation of the methylome from youthful set points, integrating stochastic CpG drift and damage-accelerated remodeling.

Ten signaling intermediates (AMPK activity, mTORC1 activity, autophagy flux, SASP output, SIRT1 activity, FOXO3 activity, PGC-1$\alpha$ activity, NF-$\kappa$B activity, p53 activity, and insulin/IGF-1 tone) are modeled as quasi-steady-state variables---algebraic functions of the primary state variables at each time step. Each is a Hill or Michaelis--Menten function of one or two state variables; for example, SIRT1 activity is $f_{\text{SIRT1}}(N) = N^{n_S} / (K_S^{n_S} + N^{n_S})$, coupling NAD$^+$ levels to downstream deacetylation. They serve as the entry points for compound effects and as the coupling terms between subsystems. Full algebraic expressions are given in Supplementary Table~S1.

\subsubsection{Ordinary differential equations}

Heteroplasmy dynamics follow the Kowald--Kirkwood replication advantage model, modified for continuous time:
\begin{equation}
\frac{dH}{dt} = \alpha \cdot H \cdot (1 - H) - \beta \cdot H \cdot \Phi_{\text{mito}}(t)
\end{equation}
where $\alpha$ is the mutant mtDNA replication advantage, $(1 - H)$ is the available wild-type population, and $\Phi_{\text{mito}}(t)$ is mitophagy flux. Mitophagy-enhancing compounds increase $\beta$.

NAD$^+$ dynamics balance synthesis via the NAMPT salvage pathway against consumption by PARP1, CD38, and SIRT1:
\begin{equation}
\frac{dN}{dt} = k_{\text{synth}} \cdot f_{\text{NAMPT}}(\text{AMPK}) - k_{\text{PARP}} \cdot D \cdot N - k_{\text{CD38}} \cdot N - k_{\text{SIRT}} \cdot N \cdot f_{\text{SIRT1}}(N)
\end{equation}
where $f_{\text{NAMPT}}$ and $f_{\text{SIRT1}}$ are Hill functions. NMN supplementation adds directly to $k_{\text{synth}}$; CD38 inhibitors reduce $k_{\text{CD38}}$.

DNA damage accumulates from ROS---amplified by heteroplasmy and SASP---and is repaired by NAD$^+$-dependent activity:
\begin{equation}
\frac{dD}{dt} = k_{\text{ROS}} \cdot (1 + \gamma_H \cdot H + \gamma_{\text{SASP}} \cdot \text{SASP}) - k_{\text{repair}} \cdot N \cdot D
\end{equation}

Methylation drift accumulates at a baseline rate accelerated by damage and SASP, partially restrained by AMPK-mediated suppression of EZH2, which reduces recruitment of DNA methyltransferases (DNMTs) to CpG sites \citep{Vire2006}:
\begin{equation}
\frac{dM}{dt} = k_{\text{drift}} \cdot (1 + \delta_D \cdot D + \delta_{\text{SASP}} \cdot \text{SASP}) - k_{\text{EZH2}} \cdot f_{\text{AMPK}}(\text{AMPK}) \cdot M
\end{equation}

Senescent cell entry follows a damage-dependent Hill function plus a constant background term $k_{\text{repl}}$ representing replicative exhaustion (Hayflick-limit arrivals per unit time); clearance is proportional to cell burden:
\begin{equation}
\frac{dS}{dt} = k_{\text{entry}} \cdot \frac{D^n}{K_D^n + D^n} + k_{\text{repl}} - k_{\text{clear}} \cdot S - k_{\text{sen}} \cdot f_Q(t)
\end{equation}
$f_Q(t)$ represents pulsed senolytic dosing, modeled as a periodic impulse (period 14 days, decay constant $\tau = 3$ days) when a senolytic compound is active.

\subsubsection{BioAge and survival}

The BioAge composite is a weighted sum of four endpoint subsystem states (NAD$^+$ is excluded because it acts as a coupling intermediate rather than an independent aging endpoint---its effects are captured through its influence on DNA repair, methylation, and mitophagy):
\begin{equation}
\text{BioAge}(t) = w_H \cdot H(t) + w_D \cdot D(t) + w_M \cdot M(t) + w_S \cdot S(t)
\end{equation}
with equal prior weights $w_H = w_D = w_M = w_S = 0.25$. The equal-prior assignment reflects agnosticism about which subsystem dominates mouse mortality and avoids consuming degrees of freedom on weight tuning. BioAge maps to instantaneous mortality hazard through a Gompertz function $h(t) = h_0 \cdot \exp(\lambda \cdot \text{BioAge}(t))$, calibrated to the UM-HET3 control survival curve. Predicted lifespan extension is the percent difference in median survival time ($T_{50}$) between treated and control simulations.

\subsection{Subsystem Weight Justification}

The equal-prior choice is deliberately simple: absent a principled empirical basis for assigning different weights, assigning equal importance to each subsystem avoids overfitting to the calibration data. As a consistency check, we mapped the component variables of GrimAge \citep{Lu2019}, PhenoAge \citep{Levine2018}, and DunedinPACE \citep{Belsky2022} to ELM-equivalent categories. The resulting weight ranges are wide (5--30\% for methylation, 20--30\% for heteroplasmy \citep{Shoubridge2007}, similar for DNA damage and senescence \citep{Ferrucci2020})---too wide to discriminate among plausible weight vectors, but narrow enough that 0.25 per subsystem is not excluded by any of them. The equal-prior choice is therefore not contradicted by external data, though it is not uniquely supported either. Sensitivity analysis (Section~3.6) establishes the biologically plausible range for each weight.

\subsection{Compound Implementation}

Each compound is implemented as a perturbation to one or more rate constants, with the injection point determined by its established mechanism of action and the magnitude derived from published dose-response data or pharmacological literature.

\begin{table}[ht]
\centering
\caption{ELM compound implementation: mechanism and injection point for all 10 modeled compounds.}
\label{tab:compounds}
\begin{tabular}{@{}lll@{}}
\toprule
Compound & Primary mechanism & ELM injection point \\
\midrule
Rapamycin & mTORC1 inhibition & $k_{\text{mTOR}}$ reduced; autophagy flux increased \\
Acarbose & $\alpha$-glucosidase inhibition $\to$ energy stress & AMPK increased; mTORC1 reduced \\
17$\alpha$-Estradiol & ER signaling; 5$\alpha$-reductase substrate & $k_{\text{drift}}$ reduced; zero effect in females \\
Canagliflozin & SGLT2 inhibition $\to$ energy stress & AMPK increased (sex-attenuated) \\
Aspirin & COX-1/2 inhibition; antioxidant & $k_{\text{ROS}}$ reduced; female: PGI$_2$ offset \\
Glycine & mTOR inhibition; antioxidant & $k_{\text{mTOR}}$ reduced (partial); $k_{\text{ROS}}$ reduced \\
Urolithin A & PINK1/Parkin mitophagy activation & $\Phi_{\text{mito}}$ increased; $H$ clearance increased \\
NMN & NAD$^+$ precursor & $k_{\text{synth}}$ increased \\
CD38 inhibitor & NAD$^+$ hydrolase inhibition & $k_{\text{CD38}}$ reduced \\
Dasatinib+Q & Senolytic: apoptosis in senescent cells & Pulsed $k_{\text{sen}}$ impulse; $f_Q(t)$ activated \\
\bottomrule
\end{tabular}
\end{table}

Dose-response curves follow a normalized Hill function $f(d) = d^n / (K_m^n + d^n)$, where $K_m$ and $n$ are fitted to ITP multi-dose data where available and estimated from literature otherwise. At 1$\times$ ITP dose, effective pathway activation ranges from $f = 0.41$ (glycine) to $f = 0.78$ (rapamycin), normalized so that $f(1) = 1$ at the calibrated ITP dose; $f > 1$ at higher doses indicates additional activation subject to diminishing returns.

\subsubsection{Sex-specific implementation}

Four sex-specific effects are implemented from independent literature, with parameter values fixed before examining ITP sex data.

(i)~\textbf{AMPK saturation:} female AMPK baseline is set to $b = 1.3\times$ male, based on estrogen-mediated upregulation of the upstream kinase LKB1 \citep{Boudaba2018,McInnes2011}. The diminishing-returns functional form is $f_{\text{atten}} = 1/(1 + k \cdot (b - 1))$, where $k$ is a global curvature parameter governing the steepness of saturation. With $k = 8$ (chosen as the midpoint of the plausible range for hyperbolic saturation of the LKB1$\to$AMPK$\to$TSC2 cascade; this is a modeling assumption, not a directly measured quantity), the attenuation factor is $1/(1 + 8 \times 0.3) = 1/3.4 \approx 0.29$. This single parameter governs the female response for all three AMPK-activating compounds.

(ii)~\textbf{CYP3A pharmacokinetics:} female effective rapamycin dose is set to 3.48$\times$ male due to reduced hepatic CYP3A clearance \citep{Lamming2014}. (iii)~\textbf{17$\alpha$-Estradiol testosterone dependency:} the compound's perturbation is set to zero in female simulations, as the 5$\alpha$-reductase mechanism requires testosterone substrate \citep{Garratt2018}. (iv)~\textbf{Aspirin COX-2 interference:} a negative PGI$_2$ offset is applied to the female aspirin perturbation \citep{Egan2004}, approximately cancelling the antioxidant benefit and producing near-zero net female effect.

\subsection{Calibration}

The calibration procedure has 6 calibrated parameters: one uniform scale factor per ITP compound, found by root-finding against male ITP targets. Approximately 190 additional parameters (ODE rate constants, BioAge weights, sex mechanism parameters) are set from published experimental measurements and structural constraints prior to calibration.

BioAge weights are equal priors ($w_M = w_D = w_H = w_S = 0.25$; justification in Section~2.2). Normalization constants are derived from control simulation values at $t = 1.0$ (baseline lifespan), ensuring $\text{BioAge}(t = 1.0) = 1.0$ by construction.

For each of the six ITP compounds, a single uniform scale factor $\alpha_c$ multiplies all of that compound's pathway injection magnitudes simultaneously (e.g., for a compound that perturbs both AMPK and mTORC1, both perturbations are scaled by the same $\alpha_c$):
\begin{equation}
\mathbf{p}_c = \alpha_c \cdot \mathbf{p}_c^{(\text{base})}
\end{equation}
where $\mathbf{p}_c^{(\text{base})}$ is the vector of pathway injection ratios for compound $c$. The scale factor is found by Brent's root-finding method (\texttt{scipy.optimize.brentq}) to satisfy $\Delta L_{\text{pred}}^{(c, \text{male})} = \Delta L_{\text{ITP}}^{(c, \text{male})}$. Male predictions are therefore exact by construction (error $< 0.01$ pp).

Female predictions are computed using the same scale factors with sex-specific mechanisms applied. No parameters are adjusted to match female targets. The six female extensions constitute genuine predictions of the calibrated model.

ITP lifespan extension values were taken from the most recent published cohort for each compound: rapamycin \citep{Harrison2009,Strong2022}, acarbose \citep{Harrison2014}, 17$\alpha$-estradiol \citep{Garratt2018}, canagliflozin \citep{Miller2020}, aspirin \citep{Strong2016}, glycine \citep{Miller2019glycine}. Where multiple cohorts tested the same compound, earlier cohorts were used for consistency checking only.

\subsection{Identifiability Analysis}
\label{sec:identifiability}

The sensitivity Jacobian $\mathbf{J}_{\text{full}} \in \mathbb{R}^{12 \times N}$ ($N \approx 196$) was constructed by finite differences (step $10^{-4}$, stability confirmed against $10^{-5}$). SVD reveals which parameter combinations are observable from ITP lifespan data and which lie in the null space. Full details in Supplementary Material.

\subsection{Sensitivity Analysis}

\textbf{Weight sweeps.} Each BioAge weight was varied from its equal-prior value of 0.25 across $[0.10, 0.50]$ in steps of 0.05, with the remaining weight ($1 - w_i$) divided equally among the other three components (since all start at 0.25, proportional and equal redistribution coincide). At each value the model was recalibrated and maximum ITP error (male calibration + female prediction) was recorded. Calibration-acceptable range: maximum ITP error $\leq$ 1.0~pp.

\textbf{Parameter perturbation.} Each of the 6 calibration parameters (compound scale factors) was independently perturbed $\pm$20\% from its calibrated value without recalibration, and the change in each key prediction was recorded.

\textbf{Gauge symmetry test.} A concern with any ODE model is that a state variable might act as a rescalable background---contributing to mortality without mechanistically differentiating among compounds. To test this for heteroplasmy, its replication advantage $\alpha$ was increased by 10\% and the fractional change in each of the 12 ITP predictions was recorded. If all predictions shift by the same fraction, heteroplasmy is merely a rescalable offset; if different compounds show different fractional changes, it is doing genuine mechanistic work.

\subsection{Combination Simulations}

Combination simulations activate multiple compound perturbations simultaneously, applied independently to their respective target rate constants. Cross-compound interactions arise only through shared state variables and signaling intermediates. Pairwise synergy is quantified as:
\begin{equation}
\Delta_{\text{syn}}(i,j) = \Delta L(i+j) - \Delta L(i) - \Delta L(j)
\end{equation}
Positive $\Delta_{\text{syn}}$ indicates super-additivity; negative indicates sub-additivity. Conditional synergy for novel targets layered on an active ITP base was computed with the base stack running in all three terms. Dose-response combinations were enumerated at $\{0, 0.5\times, 1\times, 2\times\}$ per compound.

\subsection{Software}

ODEs were integrated with SciPy's \texttt{solve\_ivp} (RK45; rtol = $10^{-8}$, atol = $10^{-10}$). Compound scale factors were calibrated by SciPy \texttt{brentq} root-finding. SVD used NumPy \texttt{linalg.svd}. Network analysis used NetworkX~3.2. Figures were generated with Matplotlib~3.8 and Plotly~5.18. All code, calibrated parameter sets, and figure generation scripts are available at \url{https://github.com/silmonbiggs/extensible-longevity-model} (MIT license).

% ============================================================
\section{Results}
% ============================================================

\subsection{Sex Differences: Four Mechanisms, Zero Sex-Specific Fitting}

The most conspicuous pattern in 20 years of ITP data is that most compounds work better in males, except rapamycin which works better in females \citep{Jiang2025}. The ELM predicts all six female outcomes from male calibration data, using four independently motivated biological mechanisms, with 1.3 percentage point mean error and zero sex-specific fitting (Table~\ref{tab:calibration}). This is the primary result of the paper.

The procedure: the model was calibrated to 6 male ITP targets by root-finding (Section~2.4), with one scale factor per compound adjusted until the male prediction matched the ITP observation. Male predictions are exact by construction. The 6 female extensions were then computed from the same scale factors with literature sex mechanisms applied. No parameters were adjusted to match female outcomes.

\begin{table}[ht]
\centering
\caption{ITP observed vs.\ ELM predicted lifespan extension. Male predictions are calibrated (6 free parameters); female predictions are genuine out-of-sample tests of the sex mechanism models. Mean female prediction error: 1.3~pp.}
\label{tab:calibration}
\begin{tabular}{@{}lccccl@{}}
\toprule
Compound & ITP Male & ELM Male & ITP Female & ELM Female & Sex mechanism \\
\midrule
Rapamycin      & +23.0\% & +23.0\% & +26.0\% & +26.7\% & CYP3A pharmacokinetics \\
Acarbose       & +22.0\% & +22.0\% & +5.0\%  & +7.6\%  & AMPK saturation \\
17$\alpha$-Estradiol & +19.0\% & +19.0\% & 0\%     & 0\%     & testosterone dependency \\
Canagliflozin  & +14.0\% & +14.0\% & +9.0\%  & +5.0\%  & AMPK saturation \\
Aspirin        & +8.0\%  & +8.0\%  & 0\%     & 0\%     & estrogen interference \\
Glycine        & +6.0\%  & +6.0\%  & +4.0\%  & +3.8\%  & AMPK saturation \\
\bottomrule
\end{tabular}
\end{table}

\textbf{AMPK saturation (acarbose, canagliflozin, glycine).} Applying the AMPK saturation model (Section~2.3.1; $b = 1.3$, $k = 8$), the female attenuation factor is 0.29---a 30\% higher AMPK baseline, on a steeply saturating curve, leaves only 29\% of the male drug response available.

The per-compound predictions: glycine is predicted well (+3.8\% predicted vs.\ +4.0\% observed), acarbose is overpredicted (+7.6\% vs.\ +5.0\%), and canagliflozin is underpredicted (+5.0\% vs.\ +9.0\%). The acarbose--canagliflozin error pattern is diagnostic. If both compounds acted purely through AMPK, their female attenuation should be similar. Instead they err in opposite directions. The most parsimonious explanation: canagliflozin has a second, sex-independent mechanism---likely FGF21/ketogenesis via SGLT2 inhibition---that the single-AMPK model does not capture. This is a hypothesis the model generates, not one it was designed to test.

\textbf{Pharmacokinetics (rapamycin).} Rapamycin is the one compound showing greater female efficacy (+26\% F vs.\ +23\% M). The AMPK saturation mechanism predicts the opposite. The reversal is pharmacokinetic: female mice achieve 3.48-fold higher blood rapamycin levels at equivalent dietary doses due to reduced CYP3A hepatic clearance \citep{Lamming2014}. Greater mTORC1 inhibition from higher circulating drug levels more than offsets the AMPK saturation disadvantage. The model predicts +26.7\% female, an error of 0.7~pp.

\textbf{Testosterone dependency (17$\alpha$-estradiol).} 17$\alpha$-Estradiol's mechanism requires testosterone as substrate via 5$\alpha$-reductase \citep{Garratt2018}. In females, without a testosterone substrate, the mechanism is inoperative---confirmed experimentally in castrated males, which lose efficacy entirely \citep{Garratt2018}. The model assigns zero effect in females; the ITP observes 0\%.

\textbf{Estrogen pathway interference (aspirin).} Aspirin's COX inhibition suppresses the COX-2$\to$prostacyclin (PGI$_2$) pathway through which estrogen mediates cardioprotection in females \citep{Egan2004,Grosser2006}. The net effect in females neutralizes a protective pathway, offsetting aspirin's antioxidant and anti-inflammatory benefits. The model predicts zero female effect; the ITP observes 0\%.

% Figure 1 placeholder
% \begin{figure}[ht]
% \centering
% \includegraphics[width=0.8\textwidth]{fig1_sex_differences.pdf}
% \caption{Sex difference predictions vs.\ ITP observations for all 6 compounds. AMPK saturation curve inset showing male vs.\ female operating points.}
% \label{fig:sex_diff}
% \end{figure}

No sex-specific parameters were adjusted to ITP lifespan data: the AMPK baseline ratio (1.3$\times$), the diminishing-returns parameter ($k = 8$), and the rapamycin PK ratio (3.48$\times$) were all set from published measurements before examining female ITP outcomes.

\subsection{Out-of-Sample Validation: Rapamycin + Acarbose}

The ITP subsequently tested rapamycin and acarbose in combination, reporting +34\% median lifespan extension in males \citep{Strong2022}. This combination was not used in any parameter fitting: neither the compound scale factors (calibrated to individual male targets) nor the BioAge weights (equal priors, fixed before calibration) were adjusted to match this result. The ELM's prediction: +33.5\%, an error of 0.5 percentage points.

The prediction is not trivially correct. Both compounds activate AMPK---rapamycin indirectly (mTORC1 inhibition relieves negative feedback on AMPK); acarbose directly (energy stress from reduced carbohydrate absorption)---and both converge on the same downstream autophagy pathway. The autophagy term follows Michaelis--Menten kinetics with $K_a = 0.35$. At ITP doses, rapamycin alone drives autophagy to $\sim$78\% of maximum activation; acarbose drives it to $\sim$65\%. Combined, the pathway encounters diminishing returns near saturation.

A naive additive model predicts +45\%. The ELM predicts +33.5\%. The ITP observes +34\%. The naive model is not close. The model identifies the direction of the interaction (sub-additive), its approximate magnitude, and the correct mechanistic reason---pathway saturation at a shared node.

% Figure 2 placeholder
% \begin{figure}[ht]
% \centering
% \includegraphics[width=0.8\textwidth]{fig2_rapa_acarb_validation.pdf}
% \caption{Rapa + Acarb validation: observed vs.\ predicted lifespan extension, naive additive shown for comparison. Autophagy Hill-function saturation curve inset.}
% \label{fig:rapa_acarb}
% \end{figure}

\subsection{Identifiability Analysis}
\label{sec:results_identifiability}

The ELM has 6 calibration parameters and $\sim$196 total parameters. The calibration is trivially well-posed: 6 parameters determined by 6 male targets via root-finding. SVD of the full sensitivity Jacobian $\mathbf{J}_{\text{full}} \in \mathbb{R}^{12 \times 196}$ reveals that the vast majority of parameters lie in the null space---set by published biology, not determined by ITP outcomes. The 6 compound scale factors are maximally constrained; predictions that depend on poorly constrained null-space parameters (SASP coupling, senolytic pulse dynamics) carry correspondingly greater uncertainty. Full identifiability analysis is provided in the Supplementary Material.

\subsection{Sensitivity Analysis and Robustness}

\textbf{Methylation weight ($w_M$).} The methylation weight ($w_M = 0.25$) was examined across $[0.10, 0.50]$ while redistributing remaining weight equally among the other three subsystems. The combination prediction ranges from $+31.8\%$ ($w_M = 0.10$) to $+37.3\%$ ($w_M = 0.50$); the female mean prediction error ranges from 1.20 to 1.56~pp. Within the biologically plausible clock-decomposition range ($w_M \in [0.10, 0.30]$), the combination prediction spans $+31.8\%$ to $+34.3\%$ and the female mean error spans 1.20 to 1.36~pp---the model is insensitive to this weight across its plausible range.

External anchoring comes from published biological age composites. Mapping the components of GrimAge2 \citep{Lu2019}, PhenoAge \citep{Levine2018}, and DunedinPACE \citep{Belsky2022} to ELM-equivalent aging categories yields a methylation fraction range of 5--30\% across the three clocks. The ELM's equal-prior $w_M = 0.25$ sits within this externally derived range.

% Figure 4 placeholder

\textbf{Heteroplasmy weight ($w_H$).} No ITP compound directly targets mitochondrial heteroplasmy, meaning $w_H$ is constrained by model structure and independent literature rather than lifespan outcomes. With equal priors ($w_H = 0.25$), the model predicts heteroplasmy accounts for 43.5\% of residual BioAge at simulated treated-animal death (vs.\ 27.9\% in simulated untreated controls). The qualitative conclusion---heteroplasmy is the primary unaddressed bottleneck among ITP compounds---holds across the full biologically plausible range of $w_H$.

% Figure 5 placeholder

\textbf{Rapa+NMN superadditivity.} The model predicts rapamycin + NMN produces superadditive lifespan extension: $+31.7\%$ combined vs.\ $+26.1\%$ additive expectation ($+23.0\%$ rapa $+$ $+3.2\%$ NMN), a superadditive component of $+5.6$~pp. Perturbing each compound scale factor $\pm$20\% individually shifts the predicted superadditive component by at most $\pm$0.9~pp; it never drops below zero across any perturbation tested.

\textbf{Gauge symmetry test.} Increasing the heteroplasmy transcription advantage by 10\% shifts all 12 ITP predictions, but different compounds shift by different amounts. The compound-specific drift ratios range from 0.151 (canagliflozin, most buffered) to 0.207 (acarbose, least buffered)---a 37\% spread---confirming that different compounds interact differently with the heteroplasmy subsystem. If the model were merely fitting a rescalable background aging rate, all drift ratios would be identical.

\subsection{Combination Predictions: ITP Compounds}

With calibration and validation established, the model was used to predict all pairwise combinations of the six ITP compounds at four dose levels ($0, 0.5\times, 1\times, 2\times$), yielding a $24 \times 24$ dose-response matrix. Dose-response curves follow normalized Hill functions fitted to ITP multi-dose data where available and estimated from pharmacological literature otherwise.

The full synergy matrix reveals a dominant pattern that reflects pathway architecture: compounds converging on the same signaling node are sub-additive; compounds acting through independent pathways are super-additive.

% Figure 6 placeholder

Among ITP pairs: rapamycin + acarbose shows sub-additive interaction (ratio 0.80), consistent with the validated Rapa+Acarb prediction and its mechanistic explanation (shared AMPK/mTORC1 pathway saturation). Rapamycin + NMN shows super-additive synergy ($+5.6$~pp above additive, ratio 1.21)---the largest positive signal among ITP compounds---because mTORC1 inhibition and NAD$^+$ replenishment address independent rate-limiting steps with no shared bottleneck.

\textbf{Dose-response analysis.} At 2$\times$ dose, effective pathway activation ranges from $f(2) = 1.02$ for acarbose (fully saturated at 1$\times$) to $f(2) = 1.50$ for glycine (still climbing). Dose escalation gains at 2$\times$ range from $+6.6$~pp (glycine) to $+14.2$~pp (acarbose), but these gains carry increasing risk and operate deep in the diminishing-returns regime. Additional extension requires novel mechanisms, not higher doses.

% Figure 7 placeholder

\textbf{Risk-benefit optimization.} Enumerating all 729 combinations at three dose levels ($0, 1\times, 2\times$) with literature-grounded risk scores produces a Pareto frontier of risk-adjusted strategies. The best ITP-only cocktail (four unisex compounds at 1$\times$) achieves $+41.1\%$ M and $+35.3\%$ F. Female optimization is dominated by rapamycin's pharmacokinetic advantage across all risk levels.

% Figure 8 placeholder

\subsection{The Heteroplasmy Bottleneck}

The BioAge composite tracks four aging subsystems simultaneously. Examining its decomposition at the moment each simulated animal reaches the death threshold---for untreated controls and for animals treated with the best unisex ITP cocktail---reveals which subsystems the model predicts are addressed by intervention and which are not.

In simulated untreated controls, BioAge at death reflects the assigned weights modified by coupled dynamics: methylation 28.7\%, heteroplasmy 27.9\%, senescence 27.5\%, DNA damage 15.9\%. The model predicts that in animals treated with the four-compound unisex ITP cocktail (rapamycin + acarbose + canagliflozin + glycine at 1$\times$), methylation, DNA damage, and senescent cell burden are substantially reduced. Heteroplasmy is not. None of the four compounds directly targets mitochondrial mutant load, and heteroplasmy's indirect connections to the AMPK/mTORC1 pathway are too weak to produce meaningful reduction through that route. A mouse receiving every available intervention still dies---of the one thing none of them addressed. At simulated treated-animal death, heteroplasmy accounts for 43.5\% of residual BioAge (vs.\ 23.1\% for methylation, 17.0\% for senescence, 16.4\% for DNA damage), becoming the dominant contributor to eventual mortality---as reliability theory predicts for a series system in which one cause of death is left untreated.

% Figure 9 placeholder

This finding must be interpreted with appropriate caution. The conclusion that heteroplasmy is the primary untreated bottleneck follows substantially from the assignment $w_H = 0.25$---a prior constrained by model structure and independent literature, not directly by ITP lifespan outcomes. The weight sensitivity analysis establishes that the qualitative conclusion holds across the full biologically plausible range of $w_H$. The quantitative share of heteroplasmy at treated-animal death is a calibrated inference, not a direct experimental measurement. It motivates the novel target predictions that follow and generates a falsifiable hypothesis: interventions that directly reduce heteroplasmy should produce the largest marginal lifespan gains when added to the ITP base.

\subsection{Novel Target Predictions}

Four targets not represented in the ITP compound set were incorporated using literature-estimated pathway injection parameters: urolithin A (mitophagy), NMN (NAD$^+$ precursor), a CD38 inhibitor (NAD$^+$ preservation via reduced consumption), and dasatinib + quercetin (senolytic, pulsed biweekly following Baker et al.\ \citep{Baker2016}). None carries ITP-calibrated dose-response data; predicted magnitudes are directional estimates. Relative ranking is more reliable than absolute magnitude for any individual compound.

Layered on the rapamycin + acarbose base (+33.5\% M), the model predicts marginal additions: NMN $+11.4$~pp, CD38 inhibitor $+7.0$~pp, dasatinib+quercetin $+0.6$~pp, urolithin A $+0.2$~pp. NMN leads because it addresses the NAD$^+$ deficit that amplifies all four aging axes simultaneously; CD38 inhibition achieves a similar effect from the consumption side. Mitophagy-targeting compounds show smaller marginals on this base because the rapa+acarb combination already partially addresses heteroplasmy through indirect AMPK$\to$FOXO3$\to$mitophagy signaling.

% Figure 10 placeholder

The primacy of mitophagy reflects the heteroplasmy bottleneck finding: urolithin A's mechanism---PINK1/Parkin and PGC-1$\alpha$ pathway activation \citep{Luan2021}---directly addresses the one aging subsystem that no ITP compound substantially reduces. Its predicted marginal addition reflects where the bottleneck is, not an assumption of extraordinary intrinsic potency.

\textbf{Conditional pairwise synergy.} The strongest cross-axis synergy is mitophagy + NMN. The mechanism is mutually reinforcing---mitophagy reduces heteroplasmy, improving OXPHOS efficiency and reducing mitochondrial drain on NAD$^+$; NMN boosts the now-more-efficient NAD$^+$ synthesis pathway. Each compound makes the other more effective.

Mitophagy + senolytic shows conditional synergy: mitophagy addresses the upstream root cause (mitochondrial mutant load), senolytics address the downstream consequence (accumulated senescent cells). This pair carries greater uncertainty than mitophagy + NMN because the senolytic predictions depend on SASP coupling parameters that are literature-estimated and not calibrated to lifespan data. NMN + CD38 inhibitor shows sub-additive conditional interaction: both target NAD$^+$ homeostasis, saturating SIRT1 activation on top of a base that already activates AMPK$\to$NAMPT.

% Figure 11 placeholder

\subsection{Full-Stack Prediction}

Simulating the four unisex ITP compounds plus NMN, CD38 inhibitor, dasatinib+quercetin, and urolithin A simultaneously at calibrated doses predicts a full-stack extension of $+69\%$ for both sexes. An internal variable audit (Supplementary Material) confirms that no model variable exceeds its single-compound calibration range by more than 35\%: the large extension arises from addressing six partially independent aging axes in parallel, not from any single variable driven to an implausible extreme.

All 28 pairwise combinations of the eight compounds show a mean interaction ratio of 1.04 (range 0.79--1.33): same-pathway pairs show diminishing returns; cross-pathway pairs show mild synergy. The full-stack prediction is a sum of many small, well-behaved interactions, not a product of combinatorial explosion. The central unvalidated assumption is whether these aging axes compound as independently in vivo as the model assumes.

Network analysis of the ODE coupling structure---node degree, betweenness centrality, and intervention path distances---provides structural explanations for the synergy patterns above. The key insight is that compounds with long path distances to BioAge (AMPK activators, distance 6.2) show the strongest interactions because their effects cascade through multiple intermediaries, while short-path compounds (methylation-targeting, distance 2.9) show predictable, low-interaction effects. Full network analysis is provided in the Supplementary Material.


% ============================================================
\section{Discussion}
% ============================================================

\subsection{The Sex Difference Explanation}

The strength of the sex-difference result lies not in the error magnitude but in the mechanistic diversity: four different biological mechanisms---a saturation curve, a pharmacokinetic ratio, a substrate dependency, and a pathway interference---each drawn from independent experimental literature, were specified before examining any female ITP data. That all four produce correct directional predictions, and four of six within 1~pp of the observed value, is evidence that the model is picking up real effects, not fitting noise.

The canagliflozin error pattern noted in Section~3.1---acarbose overpredicted, canagliflozin underpredicted, same pathway---suggests a specific experimental test. Measuring FGF21 levels in canagliflozin-treated male and female mice would determine whether the compound has a sex-independent pharmacological component via SGLT2-mediated ketogenesis. A sex-independent FGF21 response would confirm the multi-pathway hypothesis and indicate that adding a second pathway to the canagliflozin implementation would resolve the prediction error.

\subsection{Interpreting the Combination Validation}

The rapamycin + acarbose out-of-sample prediction (+33.5\% vs.\ +34\% observed) provides supporting evidence. Its value lies in what the model had to get right: predicting sub-additivity rather than the naively expected +45\% sum, for the correct mechanistic reason (pathway saturation at a shared AMPK/autophagy node). But one successful combination prediction does not validate the full model. The Rapa+Acarb interaction operates through the most tightly constrained element of the ODE system; the harder test---cross-pathway interactions involving mechanistically independent axes---requires multi-compound experimental designs that have not yet been conducted.

\subsection{Identifiability and What the Data Constrains}

The ELM's calibration is deliberately minimal: 6 compound scale factors determined by 6 male targets. The female extensions are reserved as genuine predictions rather than fitted, accepting higher error in exchange for a clean test of the sex mechanism hypotheses.

The sensitivity Jacobian SVD analysis reveals which of the $\sim$196 total parameters are observable from ITP lifespan data. The vast majority lie in the null space---set by published biology or structural constraints, not determined by ITP outcomes. This has practical consequences for interpretation: predictions that depend on the top singular vectors (heteroplasmy transcription dynamics and mTOR/AMPK signaling) deserve the most trust; predictions that depend on SASP coupling parameters---which sit near the bottom of the observable spectrum---carry greater uncertainty, as flagged in the discussion of the mitophagy + senolytic conditional synergy.

A model with 6 constrained parameters and 190 in the null space is not, in practice, a 196-parameter model. It is a 6-parameter model that has imported a library of biological literature as frozen assumptions. The model predicts correctly not despite its enormous null space but because of it: the null space is where the priors live, and those priors, being drawn from competent experimental biology, are mostly right.

But the priors are not all equally right. The edges of that imported knowledge are where the model will fail first. Predictions most sensitive to null-space parameters---SASP coupling, senolytic pulse dynamics, cancer module interactions---carry correspondingly greater uncertainty and should be interpreted as directional hypotheses rather than quantitative forecasts.

\subsection{The Heteroplasmy Bottleneck}

The prediction that mitochondrial heteroplasmy is the dominant unaddressed bottleneck in the ITP compound set is the most novel result of the paper and requires the most careful accounting of what is data and what is assumption. The inference is prior-dependent: because no ITP compound directly targets heteroplasmy, $w_H$ is set by the equal-prior assumption ($w_H = 0.25$), not by lifespan outcomes. The equal-prior choice is deliberately agnostic---it assigns no subsystem more importance than any other before seeing the data. The sensitivity analysis (Section 3.4) establishes that the qualitative conclusion holds across the entire $w_M \in [0.10, 0.50]$ sweep, with heteroplasmy consistently emerging as the dominant unaddressed contributor (43.5\% of residual BioAge at simulated treated-animal death).

Fix everything but one thing and the one thing kills you---a consequence reliability theory has formalized \citep{Gavrilov2001} and that clinical experience confirms independently of any model. The ITP compounds have substantially addressed the AMPK/mTOR, glucose metabolism, and inflammatory axes; heteroplasmy is the axis they have not touched. This is a specific, falsifiable hypothesis: direct measurement of mitochondrial heteroplasmy in post-mortem tissue from ITP-treated vs.\ control animals would confirm or refute it independently of any lifespan endpoint.

The practical implication follows from the prediction. Urolithin A (mitophagy enhancement via PINK1/Parkin and PGC-1$\alpha$ \citep{Luan2021}) is predicted to produce the largest marginal addition to the ITP base precisely because it addresses the one subsystem the base compounds leave untouched. The superadditive interaction between mitophagy and NMN arises from mutual reinforcement: mitophagy reduces heteroplasmy and improves OXPHOS efficiency, reducing the mitochondrial drain on NAD$^+$; NMN boosts the now-more-efficient synthesis pathway. This cross-axis mechanism is qualitatively different from the AMPK convergence that makes Rapa+Acarb sub-additive, and it identifies mitophagy + NAD$^+$ as the highest-priority combination to test experimentally.

\subsection{Limitations}

\textbf{Cross-pathway independence.} The central unvalidated assumption is that the four aging axes compound approximately independently in living animals. The ODE coupling structure encodes known cross-pathway interactions---the SASP$\to$NAD$^+$ drain through CD38, the ROS$\to$DNA damage feedback, the AMPK$\to$EZH2 methylation connection---and the network analysis makes these interactions explicit. But the model assumes that the interactions it encodes are the dominant ones and that higher-order feedback operating at timescales or through mechanisms not captured in the ODE system are not large enough to invalidate the combination predictions. Testing this assumption requires multi-compound lifespan experiments that have not yet been conducted.

\textbf{SASP coupling parameters.} Predictions that depend heavily on SASP dynamics---particularly the mitophagy + senolytic conditional synergy---rest on coupling parameters estimated from literature rather than calibrated to lifespan data. SASP is a heterogeneous secretory phenotype that varies substantially by cell type, tissue, and organism; the ELM's single SASP node is a deliberate simplification. These predictions should be treated as directional hypotheses rather than quantitative forecasts.

\textbf{Causes of death not modeled.} The four BioAge subsystems account for a large but not exhaustively characterized fraction of what kills these mice. Cardiovascular disease, neurodegeneration, and musculoskeletal deterioration each contribute to late-life death through pathways that overlap with but are not identical to the four modeled subsystems. A systematic census of the fraction of UM-HET3 mortality attributable to each pathway is listed as pending; the model's predictions are most reliable for outcomes where the four subsystems dominate.

\textbf{AGE crosslink dynamics.} Advanced glycation end-product crosslink accumulation was not included as a distinct state variable. The feedback routes are real: RAGE receptor activation upregulates NF-$\kappa$B, contributing to SASP and DNA damage via NADPH oxidase-derived ROS \citep{Soldatos2006}. However, on the $\sim$800-day UM-HET3 lifespan, AGE accumulation is modest compared to what occurs over decades in humans---collagen and elastin crosslinking accumulates meaningfully only on very long timescales with very long-lived proteins. On mouse timescales, the RAGE$\to$SASP contribution is likely small relative to mitochondria-driven and senescent-cell-driven SASP. More critically, no ITP-tested compound acts primarily through crosslink-breaking; an AGE crosslink ODE would therefore sit in the null space of the identifiability analysis, constrained by no calibration signal. The appropriate time to add this term is when a crosslink-breaker enters the experimental record.

\textbf{Human translation.} The model was calibrated to mouse ITP data. Translation to human biology requires pharmacokinetic scaling, tissue distribution characterization, and accounting for the much longer timescales over which AGE crosslinking, telomere shortening, and clonal hematopoiesis contribute to human mortality. The current model makes no quantitative claims about human outcomes; it provides a ranked intervention hypothesis for further testing in vertebrate models.

\subsection{Relationship to Existing Computational Aging Models}

The primary novelty of the ELM is not any individual ODE---each has antecedents in reliability theory \citep{Gavrilov2001}, mechanistic aging models \citep{Kowald1993,Kogan2015,Podolskiy2015}, and the hallmarks framework \citep{LopezOtin2013,LopezOtin2023}. It is the combination: explicit identifiability analysis, prospective out-of-sample validation against ITP sex data, and specified falsification conditions. To our knowledge, prior models in this space have not been calibrated to ITP lifespan data or validated prospectively on unseen outcomes---which raises a question about what, exactly, they are predicting. A model that cannot be wrong is not a model; it is a metaphor.

\subsection{Falsification Conditions}

Two experimental results would require model revision rather than refinement. First, direct measurement of mitochondrial heteroplasmy in post-mortem tissue from ITP-treated animals would test the bottleneck prediction independently of any lifespan endpoint. If heteroplasmy does not rise to dominate residual BioAge under treatment, the novel target rankings would change accordingly. Second, systematic measurement of epigenetic clock scores (GrimAge2 or DunedinPACE) in ITP compound cohorts would provide a partial identifiability audit from biological data. If clock-predicted trajectories diverge substantially from the methylation ODE's outputs, the methylation drift parameters would require re-examination.

\subsection{Conclusions}

The ELM demonstrates that a mechanistically structured ODE model with 6 calibrated parameters can explain the most conspicuous unexplained pattern in 20 years of ITP data: why most longevity compounds work better in males, except rapamycin which works better in females. Four independently motivated sex mechanisms---AMPK/LKB1 saturation, CYP3A pharmacokinetics, testosterone dependency, and COX-2/prostacyclin interference---predict all six female outcomes from male calibration data with 1.3~pp mean error and zero sex-specific fitting. The model also predicts the rapamycin + acarbose combination outcome with 0.5~pp error, for the correct mechanistic reason.

The most actionable prediction is the heteroplasmy bottleneck: that after the best available ITP compound stack, mitochondrial heteroplasmy is the dominant unaddressed contributor to residual aging, and that mitophagy-enhancing compounds are therefore predicted to produce the largest marginal lifespan addition to that base. This prediction is prior-dependent---$w_H$ is constrained by literature and model structure, not by lifespan data---and its quantitative precision should be interpreted accordingly. It generates a specific experimental hypothesis testable independently of any lifespan endpoint: direct measurement of mitochondrial heteroplasmy in post-mortem tissue from ITP-treated vs.\ control animals.

The longevity intervention field currently lacks a quantitative model for predicting combination outcomes from single-compound data. The ELM is a first step: calibrated to the best available mammalian lifespan data, validated out-of-sample on sex differences and a drug combination, and designed to inform decisions about which combinations may warrant testing---and why.

% ============================================================
% Bibliography
% ============================================================

\bibliographystyle{unsrtnat}
\bibliography{references}

\end{document}
